\chapter{The applications built on the library}
\label{ch:applications}

The library is exercised by a set of applications that resemble the components
of a trading venue: gateways, a sequencer, a matching engine, an authentication
service and a publisher.

\textbf{They are not a trading venue, and are not intended to become one.} They
exist to put the library under the kinds of load and failure a venue produces,
which is a different purpose from trading. The matching engine, for instance,
accepts and cancels orders and does not match them. Absences of that kind are
deliberate, and this chapter records which they are, so that a reader can tell
them from defects.

The rest of this chapter is not yet written, except for the order gateways' cancel
requirements, their limits on how fast a member may send, and the placeholders below.

\section{The order gateways}
\label{sec:order_gateways}

Two gateways accept orders from members. The FIX gateway speaks FIX 5.0 SP2 over a
FIXT.1.1 session. The binary gateway speaks an encoding of the venue's own. Members are
moving towards the binary gateway for its lower latency, and both are expected to remain.

\subsection{The fields a request to cancel must carry}
\label{subsec:cancel_required_fields}

FIX 5.0 SP2 requires four things on an OrderCancelRequest: ClOrdID, the instrument, Side
and TransactTime. It makes OrigClOrdID optional, because a member may instead name the
order by the OrderID the venue assigned it. It makes the quantity optional as well.

The venue identifies a member's order by the identifier the member submitted it under. The
matching engine keys its book on that identifier, and the venue offers a member no other
way to name an order it has placed. A request to cancel therefore carries it.

\begin{requirement}{R-0142}{A request to cancel an order names the order it cancels}
  A request to cancel an order shall carry the identifier the order was submitted under.
  The venue shall refuse a request that omits it.
  \rationale{The identifier on the cancel request is the request's own, not the order's.
  A request carrying only that identifier names nothing the venue can look up.

  FIX permits a member to name the order by the identifier the venue assigned instead.
  The venue does not accept that alternative, because the matching engine's book is keyed
  on the member's identifier and holds no index on its own.}
  \uncovered{no scenario sends a cancel that omits the original identifier}
\end{requirement}

The venue requires nothing else beyond what FIX requires. In particular it does not
require the order's quantity. A cancel applies to the whole of an order, so the quantity
cannot change what the venue does with the request, and a member that sends one is sending
a field the venue reads and discards.

\begin{requirement}{R-0144}{Both gateways accept the same request to cancel}
  The two gateways shall require the same fields on a request to cancel an order. Where
  FIX 5.0 SP2 makes a field optional, the venue shall accept a request that omits it,
  unless this document states a requirement for that field.
  \rationale{A member choosing between the two gateways is choosing on latency. A member
  that has to send different messages to the two of them is choosing on message format as
  well, and a member moving from one to the other has to change its code to do it.

  The requirement is written as a rule about the pair rather than as a list of fields,
  because the fields will change and the agreement between the gateways is the property
  worth keeping.}
  \uncovered{no scenario sends the same cancel to both gateways}
\end{requirement}

\subsection{What a member is told about a message the venue will not process}
\label{subsec:refused_message_reply}

A member that sends a message the venue refuses is told so. The gateways answer a refused
new order with an execution report that rejects the order, and a refused request to cancel
with a reply that refuses the request and reports the order as still open. Each answer
names the reason, such as the fields that were missing.

\begin{requirement}{R-0143}{A message missing a field the venue requires is answered}
  Where the venue refuses a member's order or request to cancel because a field the venue
  requires is absent, the venue shall answer the member, and shall name the absent fields.
  Where the message omits the identifier the answer would be addressed by, the venue may
  discard it.
  \rationale{A member whose message is discarded without a reply cannot tell refusal from
  loss. Of the two readings, the member takes the one where the message is still being
  acted on, which is the more dangerous. A member that believes a cancel was accepted, and
  whose order is in fact still resting, is exposed to a market it believes it has left.

  Naming the absent fields is what makes the answer usable. A member told only that fields
  were missing has to search a message of thirty fields for the one the venue objected to,
  and the venue already knows which it was.

  The exception is narrow. An answer is addressed by the identifier the member put on the
  message, so a message without one cannot be answered at all.}
  \uncovered{no scenario sends a message with a required field absent}
\end{requirement}

\begin{requirement}{R-0151}{A refused request to cancel is not answered as a rejected order}
  Where a gateway refuses a member's request to cancel or amend an order, it shall answer
  with a reply that refuses the request, and shall not answer with a reply that rejects the
  order. The reply shall report the order's current status. In FIX, the reply shall be an
  OrderCancelReject.
  \rationale{The order a refused request names is still on the book. FIX gives an execution
  report rejecting an order one meaning: that the order is not on the book. A member that
  receives one in answer to a cancel believes it has no order in the market, while the
  venue can still trade the order it holds. A member exposed to a market it believes it has
  left is the most dangerous outcome a reply can produce.

  The rule holds whatever the reason for refusing the request: a missing field, the venue
  not accepting orders, or a limit on how fast the member may send (\reqref{R-0150}).}
  \covers{ha_test:42}
\end{requirement}

\begin{requirement}{R-0152}{Both gateways refuse the same commands for the same reasons}
  Each gateway shall check a member's order or request to cancel before passing it to the
  sequencer, and shall refuse one that has a required field empty, an enumerated field
  holding a value the venue does not define, a quantity or price that is not a decimal
  number, or an identifier, symbol or quantity longer than the venue holds. Each gateway
  shall also refuse an order or request to cancel while no sequencer is connected or the
  venue is not accepting orders. A refusal shall be answered, and shall say why.
  \rationale{A member cannot be held to two sets of rules according to the protocol it
  happens to speak. A command passed on unchecked is worse than one refused: an identifier
  longer than the matching engine's book key is cut short there, so two orders can become
  one, and a command the engine cannot read is dropped with no reply at all.

  A binary command has to be decoded to be checked, which the FIX gateway's parsing
  already does. Decoding fields at fixed positions costs less than parsing FIX text, so
  the binary gateway keeps its advantage.}
  \covers{ha_test:57}
  \covers{ha_test:42}
\end{requirement}

\subsection{How fast a member may send}
\label{subsec:gateway_throttles}

The venue limits how fast each session may send three kinds of command: placing an order,
amending an order and cancelling an order. Each limit is a number of commands per second.
The venue provisions the limits for each comp id, and the gateways apply them to each
session separately. The design is in \texttt{docs/venue/gateway\_throttles.md}.

\begin{requirement}{R-0148}{Each session is limited in how fast it may place, amend and cancel}
  The venue shall hold, for each comp id, a limit on the number of commands per second for
  each of placing, amending and cancelling an order. Each limit shall be a whole number from
  0 to 100,000, where 0 means the commands of that kind are not limited, and 0 shall be the
  value of a limit nobody has set. The gateways shall apply a comp id's limits to each of its
  sessions separately, with the values that held when the session logged on.
  \rationale{A member sending faster than the venue can serve takes capacity from every
  other member. A limit for each comp id lets the venue give each member the rate it has
  agreed to, and nothing beyond it.

  The limits apply to each session separately because the gateways count commands where
  they arrive. Counting across all the sessions of a comp id would need the gateways to
  share their counts with each other, on every command.

  The venue processes an order from start to end in tens of microseconds, so 100,000
  commands per second, one every 10 microseconds, is the fastest rate any member could
  use. Amending is limited although the venue cannot yet amend an order, so that amending
  is limited from the day it exists.}
  \covers{ha_test:56}
\end{requirement}

\begin{requirement}{R-0149}{A limit holds over every one-second period}
  A gateway shall accept a command only if fewer commands of the same kind than the limit
  were accepted on the session during the second before it arrived. A command the gateway
  refuses shall not count towards the limit.
  \rationale{Measuring over the second before each command means no one-second period,
  wherever it starts, contains more commands than the limit. A count that started again at
  each whole second would allow twice the limit across the boundary between two seconds.

  A member that keeps sending while it is being refused would never be accepted again if
  its refused commands counted. Only accepted commands use the capacity the limit
  protects.}
  \covers{ha_test:56}
\end{requirement}

\begin{requirement}{R-0150}{A command over the limit is refused, and the member is told why}
  A gateway shall refuse a command that would exceed its session's limit, and shall answer
  the member with the reason, naming the limit. The gateway shall not pass the command to
  the sequencer, shall not hold it to send later, and shall keep the session open.
  \rationale{A member must be able to tell a refused command from a lost one, for the reason
  \reqref{R-0143} gives. A command held and sent later would reach the matching engine at a
  time the member did not choose, against a market that may have moved. Closing the session
  would cancel the member's resting orders under cancel-on-disconnect, which is a far
  heavier consequence than sending too fast.}
  \covers{ha_test:56}
\end{requirement}

\section{The matching engine publisher}
\label{sec:matching_engine_publisher}

\emph{Placeholder, created 2026-09-08. Nothing here is settled.}

The matching engine publisher follows the sequencer's write-ahead log and republishes what it
sees. It is the venue's first publisher, and today it is the only one.

\subsection{What it publishes}
\label{sec:mep_pubsub}

Two topics exist: orders, and execution reports. This subsection is to record what each carries,
who is entitled to subscribe to it, and what a subscriber is promised --- as requirements, and
referring back to section~\ref{sec:pubsub} for anything general rather than restating it.

Two facts to keep in view while writing it. The publisher stores each record under the
\emph{sequencer's} sequence number rather than one of its own, so a position means the same thing
on either instance of the pair and a subscriber's cursor survives a failover. And the fanout as
built is a loop over unicast connections, which is the scaffolding the pub/sub mechanism was
always meant to replace.

\subsection{Technical events}
\label{sec:mep_technical_events}

\emph{Placeholder, created 2026-09-08. Whether this subsection belongs to this component at all
is undecided; see below.}

\textbf{The requirements are not here.} They belong to
\cref{sec:ha:technical_event}, because the venue's inability to say anything about itself is what
blocks the requirement that it not run with no instance of a component without saying so. What
belongs in this subsection is this component's part in carrying the announcement, referring there
rather than restating it.

Which leaves the question this subsection waits on: whether this component publishes the technical
event at all. \Cref{sec:ha:technical_event} establishes two audiences reached by two paths.
A component that must act on a phase --- a gateway refusing a member's order, the sequencer
declining to sequence one --- takes it on the control plane, most likely by the existing
order-acceptance broadcast gaining the phase and the halted flag rather than by a second mechanism
competing with it. A component that only needs to know takes it by publish and subscribe, and this
publisher is named among those.

So this publisher appears there as something to be \emph{told}. Whether it is also the component
that owns the topic and redistributes the event is open. There is an argument that it is: a
topic's history lives in exactly one publisher's write-ahead log, this publisher already follows
the sequencer's log, and it stores each record under the sequencer's own sequence number --- so an
event that entered through the sequencer would keep its position relative to the orders it speaks
about, which is the whole of what a halt means. There is an argument that it is not: a component
that both consumes a fact and republishes it is doing two jobs, and the venue may prefer the event
to reach its subscribers from wherever it is declared.

That is to be settled with the trading-halt entry in the bug list, and in the order of work set
out in \texttt{docs/venue/trading\_phases.md} --- not here and not first.
